\section{Marco Teórico}
\label{sec:theory}

Esta sección desarrolla un marco estilizado basado en tareas que organiza las predicciones empíricas y aclara por qué la respuesta del mercado laboral a los LLM en América Latina es ex ante ambigua. El marco adapta a \citet{AcemogluAutor2011_tasks} y \citet{AcemogluRestrepo2018_race, AcemogluRestrepo2022_tasks_inequality} al caso de una tecnología de propósito general que reconfigura principalmente tareas cognitivas no rutinarias. El modelo es deliberadamente compacto: su función es disciplinar la especificación empírica, no entregar una calibración cuantitativa.

\subsection{Producción basada en tareas}
\label{subsec:theory_tasks}

El producto final en la ocupación $o$ se obtiene combinando un continuo de tareas $j \in [0, N]$. Cada tarea es realizada por trabajo humano o por capital, donde el capital incluye la automatización tradicional y, a partir de 2022, software construido sobre LLM. Las ocupaciones son paquetes de tareas: la ocupación $o$ se caracteriza por un vector fijo de participaciones $\{s_{oj}\}_{j \in [0,N]}$ con $\sum_j s_{oj} = 1$, donde $s_{oj}$ es la participación del producto de la ocupación $o$ generado por la tarea $j$. La tecnología que define a una ocupación se toma como fija en el corto plazo; la reasignación entre ocupaciones es regida por la demanda laboral.

El salario en la ocupación $o$ es la suma de la remuneración específica por tarea:
\begin{equation}
w_o \;=\; \sum_j s_{oj} \cdot p_j,
\label{eq:wage_basic}
\end{equation}
donde $p_j$ es el precio de equilibrio de la tarea $j$. Antes de los LLM, $p_j$ depende de la productividad del trabajo humano en la tarea $j$ y de la demanda local por bienes que usan intensivamente la tarea $j$. El marco sigue a \citet{AutorLM2003_skill_content} al tratar la distinción entre tareas rutinarias y no rutinarias como el eje relevante para olas previas de automatización, y sigue a \citet{AcemogluAutor2011_tasks} al permitir que el capital sustituya al trabajo en cualquier tarea para la que exista una tecnología suficientemente productiva.

\subsection{Los LLM como una nueva tecnología}
\label{subsec:theory_llm}

El lanzamiento de ChatGPT en noviembre de 2022 marca la fecha en la que un subconjunto de tareas $J^{\text{LLM}} \subset [0, N]$ pasa a ser realizable por software construido sobre LLM. El puntaje de exposición complementaria de Eloundou et al. es la participación de las tareas de la ocupación $o$ que caen en este subconjunto:
\begin{equation}
\beta_o \;\equiv\; \sum_{j \in J^{\text{LLM}}} s_{oj}.
\label{eq:beta_def}
\end{equation}
El puntaje $\beta_o$ mide la participación de las tareas de una ocupación en las que al menos el 50 por ciento del tiempo requerido puede reducirse mediante software \emph{construido sobre LLM}, en contraste con la sustitución directa por LLM ($\alpha_o$) o la suma cota superior ($\zeta_o = \alpha_o + \beta_o$); véase \citet{Eloundou2023_gpt_exposure}. La elección de $\beta$ y no de $\zeta$ está motivada teóricamente para el contexto de ALC: la interacción directa con interfaces de chat sigue siendo modesta, pero el software de terceros con capacidad LLM incorporada (plataformas de relación con clientes, herramientas de contabilidad, chatbots de atención al cliente, asistentes de código) es el canal de difusión principal a través de las firmas multinacionales que operan en la región.

La llegada de los LLM tiene dos efectos opuestos sobre el salario en las ocupaciones expuestas. El \emph{efecto desplazamiento} ($\tau_d > 0$) opera porque las firmas ahora producen las tareas $j \in J^{\text{LLM}}$ a menor costo usando software, reduciendo la demanda laboral en las ocupaciones correspondientes. La magnitud escala con $\beta_o$: las ocupaciones con mayor participación de tareas automatizables enfrentan un desplazamiento más fuerte. El \emph{efecto reposición} ($\tau_r > 0$) opera porque las firmas crean nuevas tareas complementarias a la tecnología: diseño de prompts, verificación de salidas, supervisión, integración con flujos de trabajo existentes. Estas tareas son realizadas por humanos junto al software y elevan la productividad del trabajo en las ocupaciones expuestas.

La respuesta salarial neta es la diferencia de los dos canales, escalada por la exposición:
\begin{equation}
\Delta w_o \;=\; \bigl(\tau_r - \tau_d\bigr) \cdot \beta_o.
\label{eq:net_effect}
\end{equation}
El signo y la magnitud de $\tau_r - \tau_d$ son objetos empíricos, no predicciones teóricas. El marco solo fija la forma entre ocupaciones: la respuesta es proporcional a $\beta_o$ en valor absoluto, y las ocupaciones de alta exposición se mueven más (en cualquier dirección) que las de baja exposición.

\subsection{Adopción heterogénea: el canal de la formalidad}
\label{subsec:theory_formality}

La adopción de software aumentado con LLM no es gratuita a nivel de la firma. Las cuotas de suscripción, la integración con sistemas existentes, la capacitación de los trabajadores y el soporte de TI requieren inversión de costo fijo. Las firmas con débil acceso al crédito, baja capitalización y escasa capacidad gerencial enfrentan costos de adopción prohibitivos. En América Latina, esta heterogeneidad se mapea estrechamente sobre la división formal-informal: las firmas formales son más grandes, están más capitalizadas y tienen mayor probabilidad de operar en sectores donde el software aumentado con LLM ya está comercializado; las firmas informales son más pequeñas, restringidas en efectivo y operan con frecuencia fuera de la infraestructura digital requerida para desplegar estas herramientas.

El marco predice, por tanto, efectos asimétricos de corto plazo a través de la división formal-informal. Dentro del sector formal operan ambos canales: $\tau_d^F > 0$ genera desplazamiento parcial y $\tau_r^F > 0$ genera creación de tareas complementarias, con magnitudes relativas que pueden favorecer a los trabajadores con educación terciaria capaces de realizar las nuevas tareas de supervisión. Dentro del sector informal, la adopción directa es cercana a cero en el corto plazo, de modo que $\tau_d^I \approx 0$ y $\tau_r^I \approx 0$ en el sentido de equilibrio parcial. Los trabajadores informales pueden, no obstante, verse afectados indirectamente: las firmas formales que adoptan la tecnología se vuelven más productivas y pueden capturar cuota de mercado a expensas de competidores informales, generando un desplazamiento indirecto por el lado de la demanda sobre el empleo informal cuyo signo es difícil de fijar a priori.

La comparación neta es:
\begin{equation}
\Delta w_o^F - \Delta w_o^I \;=\; \bigl[(\tau_r^F - \tau_d^F) - (\tau_r^I - \tau_d^I)\bigr] \cdot \beta_o.
\label{eq:formal_informal}
\end{equation}
Que el diferencial formal-informal sea positivo o negativo depende de la fuerza relativa de la adopción directa en las firmas formales (que puede ir en cualquier dirección) y de los derrames indirectos sobre las firmas informales (típicamente negativos). El signo es una pregunta empírica en cada país.

\subsection{Predicciones contrastables}
\label{subsec:theory_predictions}

El marco entrega cuatro predicciones que la especificación empírica de la \cref{sec:method} está diseñada para evaluar.

\textbf{P1 (entre ocupaciones, dentro del país).} La respuesta del mercado laboral en el país $c$ es proporcional a $\beta_o$:
\begin{equation*}
\frac{\partial Y_{ot}^c}{\partial \beta_o} \;\gtreqless\; 0
\end{equation*}
según si la reposición domina al desplazamiento, lo iguala o es dominada. El signo es empírico y puede diferir entre países.

\textbf{P2 (formal vs. informal, dentro del país).} La adopción de corto plazo se concentra en el sector formal. El marco predice $|\tau^F| > |\tau^I|$ en el horizonte de dos años para los países en los que la adopción formal es relevante. En países donde la adopción formal también es baja, ambos canales quedan amortiguados.

\textbf{P3 (comparación entre países de las estimaciones).} Los países con mayores tasas de formalidad y mejor infraestructura digital deberían exhibir un $|\tau|$ mayor que los países con menor formalidad e infraestructura más débil. Esta predicción se evalúa comparando las seis estimaciones específicas por país $\{\hat\tau^c\}_{c=1}^6$ frente a los indicadores correspondientes de formalidad y preparación digital a nivel de país. La heterogeneidad entre países es el resultado principal y no una prueba de robustez.

\textbf{P4 (heterogeneidad por habilidad dentro de la ocupación).} Dentro de una ocupación expuesta, los trabajadores con educación terciaria están mejor posicionados para realizar las nuevas tareas de reposición (supervisión, prompting, verificación) y, por lo tanto, experimentan un efecto desplazamiento menor que los trabajadores con educación secundaria en la misma ocupación. Esta predicción se evalúa interactuando $\beta_o \cdot \text{Post}_t$ con la educación del trabajador en cada país.

\subsection{Por qué ALC difiere de la evidencia de la OCDE}
\label{subsec:theory_lac}

El marco reconcilia las estimaciones aparentemente divergentes en la literatura causal existente. \citet{Brynjolfsson2023_genai} y \citet{CuiDemirerJaffe2025_copilot} documentan grandes ganancias de productividad intra-firma en contextos estadounidenses donde la adopción se entrega de manera exógena a trabajadores de alta formalidad y preparación digital. Estas estimaciones corresponden a un entorno de alto $\tau_r$ en firmas del sector formal: el canal de reposición está plenamente operativo porque las instituciones de soporte están en su lugar. \citet{HumlumVestergaard2025_llm} documentan efectos agregados cercanos a cero en Dinamarca, un contexto con alta formalidad, alta adopción y mercados laborales ajustados. El marco racionaliza este resultado como un entorno de alto $\tau_d$ y alto $\tau_r$ en el que los dos canales aproximadamente se compensan, dejando la respuesta salarial agregada cercana a cero.

ALC difiere en ambos márgenes. La adopción es heterogénea entre países (y entre estratos de formalidad dentro de los países), y los mercados laborales no son ajustados en el sentido danés. La variación entre países en formalidad (desde aproximadamente 25 por ciento informal en Uruguay hasta más del 70 por ciento en partes de la región andina) genera predicciones ex ante que varían según el país. Por eso una estrategia empírica pa\'{i}s por pa\'{i}s es la prueba correcta del marco. Agrupar promediar\'{i}a sobre una heterogeneidad que la teoría predice es de primer orden. Cada una de las seis estimaciones por país es una prueba independiente del mismo mecanismo subyacente, y el patrón entre países de las estimaciones es la contribución empírica central.
